%
% File acl2018.tex
%
%% Based on the style files for ACL-2017, with some changes, which
% were, in turn,
%% Based on the style files for ACL-2015, with some improvements
%%  taken from the NAACL-2016 style
%% Based on the style files for ACL-2014, which were, in turn,
%% based on ACL-2013, ACL-2012, ACL-2011, ACL-2010, ACL-IJCNLP-2009,
%% EACL-2009, IJCNLP-2008...
%% Based on the style files for EACL 2006 by
%%e.agirre@ehu.es or Sergi.Balari@uab.es
%% and that of ACL 08 by Joakim Nivre and Noah Smith

\documentclass[11pt,a4paper,onecolumn]{article}
\usepackage[hyperref]{acl2018}
\usepackage{times}
\usepackage{latexsym}

\usepackage{url}

%\aclfinalcopy % Uncomment this line for the final submission

%\setlength\titlebox{5cm}
% You can expand the titlebox if you need extra space
% to show all the authors. Please do not make the titlebox
% smaller than 5cm (the original size); we will check this
% in the camera-ready version and ask you to change it back.

\newcommand\BibTeX{B{\scib}\TeX}

\title{Instructions for ACL 2018 Proceedings}

\author{First Author \\
  Affiliation / Address line 1 \\
  Affiliation / Address line 2 \\
  Affiliation / Address line 3 \\
{\tt email@domain}}

\date{}

\begin{document}
\maketitle
\begin{abstract}

\end{abstract}

\section{Introduction}

\subsection{Background}
Social media platforms host a vast amount of information, including
images, text, and videos. These data sources have become valuable for
various social science research fields, such as regional political
studies and hate speech analysis. To facilitate such research, web
crawlers are often employed to extract relevant content from social
media. However, raw crawled data cannot be directly utilized; they
require preprocessing, such as image, text, or video classification,
before meaningful analysis can be conducted.

\subsection{Problem Statement}
For social science research focused on Nazi-related content, the
ability to classify images containing Nazi symbols is crucial.
However, existing commercial solutions for Nazi symbol detection and
classification are expensive, making them impractical for analyzing
large-scale datasets. Therefore, there is a pressing need for a
cost-effective and accessible solution.

An effective classification system must not only determine whether an
image contains Nazi symbols but also identify the specific type of
symbol present. Common Nazi-related symbols include the swastika, SS
skull, and Siegrune, among others. Developing a robust classification
system that accurately detects and categorizes these symbols is
essential for researchers working in this domain.

\subsection{Objectives}
This study aims to develop an efficient image classification
methodology for detecting and categorizing Nazi symbols. To achieve
this, we compiled a dataset of approximately 10,000 images containing
13 distinct Nazi-related symbols alongside non-Nazi imagery. While
this dataset represents only a small subset of real-world Nazi
symbols, our goal is to propose a generalizable and efficient
approach that can be applied beyond the current dataset.

\subsection{Scope and Limitations}

This study focuses on two primary classification tasks:

\begin{enumerate}
  \item \textbf{Binary classification:} classify whether a given
    image contains nazi symbols or not.
  \item \textbf{Multi-class classification:} find out the which nazi
    symbols presents.
\end{enumerate}

The classification system is designed to recognize the following 13
Nazi-related symbols:

\begin{itemize}
  \item black sun
  \item british union of fascist
  \item broken sun cross
  \item happy merchant
  \item hilter
  \item hilter salute
  \item judenstern
  \item neo nazi
  \item siegrune
  \item ss skull
  \item sturmabteilung emblem
  \item swastika
  \item wolfsangel
\end{itemize}

The \textbf{Neo-Nazi} class includes symbols from various far-right
extremist groups, such as National Rebirth of Poland, Combat 18,
Atomwaffen Division, Kolovrat, Volksfront, Celtic Cross, Hammerskins,
Identitäre Bewegung, Blood & Honor, and Golden Dawn. Since only a
limited number of images were available for each group, they were
consolidated under a single "Neo-Nazi" class for classification purposes.

Beyond dataset constraints, the study also faces hardware limitations:

\begin{itemize}
  \item no GPU from the social science department to use due to configurations
  \item only google colab with 4 GB GPU available.
\end{itemize}

Given these constraints, the classification system must be
lightweight and computationally efficient to ensure effective model
training and inference within the available resources.

\subsection{Thesis Structure}

This thesis is organized as follows:

\begin{itemize}
  \item \textbf{Chapter 2: Literature Review} – We examine existing
    research and methodologies related to Nazi symbol detection and
    classification, highlighting key findings and identifying gaps
    that our work aims to address.
  \item \textbf{Chapter 3: Methodology} – This chapter details our
    approach, including data collection, dataset composition,
    preprocessing techniques, and the machine learning algorithms
    used for classification. We also discuss model training,
    evaluation metrics, and the rationale behind our chosen methods.
  \item \textbf{Chapter 4: Results} – We present the outcomes of our
    experiments, providing quantitative and qualitative analyses of
    the classification system’s performance.
  \item \textbf{Chapter 5: Discussion and Analysis} – We interpret
    the results, discuss their implications, and compare our findings
    with previous research. We also explore the system's limitations
    and potential improvements.
  \item \textbf{Chapter 6: Conclusion and Future Work} – This final
    chapter summarizes our contributions, outlines the key takeaways,
    and suggests directions for future research in Nazi symbol classification.
\end{itemize}

\section{Literature Review}

\section{Methodology}
In this section, we will discuss our training data first, then we
discuss about data preprocessing steps, at the end we talk about the
algorithms we used.

\subsection{Data}

The data are collected mainly from platform Roboflow, kaggle etc.
Nazi symbols are mainly from Roboflow with labels. Originally,
there are 23 different nazi symbols. Due to the small amount of some
neo nazi symbols, we group them into the new class
neo-nazi.
[figure data distribution of all data]

For the nazi symbol detection task, we group all the nazi symbols
into the class nazi. And for the nazi symbol classification
task, we use the data except non-nazi class.
[figure data distribution of nazi and non nazi class]

\section{Results}

\section{Discussion}

\section{Ethical Considerations}

\section{Conclusion}

% include your own bib file like this:
%\bibliographystyle{acl}
%\bibliography{acl2018}

\end{document}
